

\documentclass[conference]{IEEEtran}
\IEEEoverridecommandlockouts


\usepackage[utf8]{inputenc}
\usepackage[T1]{fontenc}
\usepackage[french]{babel}
\usepackage{cite}
\usepackage{amsmath,amssymb,amsfonts}
\usepackage{algorithmic}
\usepackage{graphicx}
\usepackage{textcomp}
\usepackage{xcolor}
\usepackage{booktabs}
\usepackage{multirow}
\usepackage{array}
\usepackage{url}
\usepackage{hyperref}
\hypersetup{colorlinks=true, linkcolor=black, citecolor=blue, urlcolor=blue}


\newcommand{\RQ}[1]{\textbf{RQ#1}}
\newcommand{\H}[1]{\textbf{H#1}}

\def\BibTeX{{\rm B\kern-.05em{\sc i\kern-.025em b}\kern-.08em
    T\kern-.1667em\lower.7ex\hbox{E}\kern-.125emX}}

\begin{document}



\title{
AdaptiveHybrid : Système de Recommandation Hybride Adaptatif\\
combinant KNN-CF et SmartCBF pour la Gestion du Démarrage à Froid\\
{\footnotesize
\textsuperscript{}Examen final --- Module Méthodologie de Recherche ---
ENSET Mohammedia, 2025--2026}
}

\author{
\IEEEauthorblockN{1\textsuperscript{ère} Hafssa MIFTAH IDRISSI}
\IEEEauthorblockA{
\textit{Département Mathématiques et Informatique} \\
\textit{ENSET Mohammedia, Université Hassan II} \\
Mohammedia, Maroc \\
hafssa.miftah.idrissi@gmail.com}
\and
\IEEEauthorblockN{2\textsuperscript{ème} Hanane AIT LHAJ}
\IEEEauthorblockA{
\textit{Département Mathématiques et Informatique} \\
\textit{ENSET Mohammedia, Université Hassan II} \\
Mohammedia, Maroc \\
hananeaitlhaj12@gmail.com}
}

\maketitle



\begin{abstract}
	% Partie Française
	Les systèmes de recommandation collaboratifs souffrent du problème du
	démarrage à froid (\textit{cold-start}) : ils sont incapables de produire des
	suggestions pertinentes pour les utilisateurs ayant peu ou pas d'historique
	de notation. À l'inverse, les approches fondées sur le contenu
	(\textit{Content-Based Filtering}, CBF) peinent à capturer les préférences
	latentes des utilisateurs actifs. Aucune de ces deux familles, prise
	isolément, ne constitue une solution universelle. Ce gap motive l'exploration
	de méthodes hybrides capables de combiner leurs forces respectives. Nous
	proposons \textit{AdaptiveHybrid}, un système hybride adaptatif combinant un
	filtre collaboratif KNN (KNN-CF) enrichi et un module de contenu amélioré
	(SmartCBF) intégrant pondération TF-IDF des genres, repli démographique et
	\textit{prior} de popularité. La pondération entre les deux composantes
	(paramètre $\alpha$) est apprise par validation croisée, et un mécanisme de
	correction du démarrage à froid ajuste dynamiquement $\alpha$ en fonction de
	l'activité de l'utilisateur. Le protocole expérimental comprend 5 exécutions
	indépendantes, trois métriques complémentaires (RMSE, NDCG@10, Recall@10) et
	deux jeux de données de référence : MovieLens 100K et MovieLens 1M. Les
	résultats montrent qu'AdaptiveHybrid surpasse la Popularité sur RMSE et
	dépasse SVD sur les métriques de classement NDCG@10 et Recall@10 sur ML-1M,
	tout en offrant une meilleure gestion du démarrage à froid grâce au repli
	SmartCBF. Ces résultats confirment nos deux hypothèses principales :
	(H1) la combinaison CF+CBF est meilleure que chacune des composantes seules,
	et (H2) cette combinaison adaptativement atténue significativement le démarrage à froid.
	
	\vspace{0.8em}
	\noindent\textit{\textbf{Abstract}---}Collaborative filtering recommendation systems suffer from the cold-start problem: they are unable to produce relevant suggestions for users with little to no rating history. Conversely, Content-Based Filtering (CBF) approaches struggle to capture the latent preferences of active users. Neither of these two families, taken in isolation, constitutes a universal solution. This gap motivates the exploration of hybrid methods capable of combining their respective strengths. We propose AdaptiveHybrid, an adaptive hybrid system combining an enriched K-Nearest Neighbors Collaborative Filtering (KNN-CF) module and an enhanced content module (SmartCBF) integrating TF-IDF genre weighting, demographic fallback, and popularity priors. The weighting between the two components (parameter $\alpha$) is learned through cross-validation, and a cold-start correction mechanism dynamically adjusts $\alpha$ based on user activity. The experimental protocol includes 5 independent runs, three complementary metrics (RMSE, NDCG@10, Recall@10), and two benchmark datasets: MovieLens 100K and MovieLens 1M. Results show that AdaptiveHybrid outperforms Popularity-based baselines on RMSE and surpasses SVD on ranking metrics (NDCG@10 and Recall@10) on ML-1M, while offering superior cold-start management through the SmartCBF fallback. These results confirm our two primary hypotheses: (H1) the CF+CBF combination outperforms each individual component, and (H2) this adaptive combination significantly mitigates the cold-start problem.
\end{abstract}

\begin{IEEEkeywords}
	Systèmes de recommandation, Filtrage collaboratif, Filtrage par contenu, Hybridation adaptative, Démarrage à froid, KNN, TF-IDF, NDCG, MovieLens, Recommender systems, Collaborative filtering, Content-based filtering, Cold-start.
\end{IEEEkeywords}

\section{Introduction}
\label{sec:intro}

Avec l'explosion du volume de contenus disponibles sur les plateformes
numériques streaming vidéo, e-commerce, musique la surcharge
informationnelle est devenue un problème central de l'expérience utilisateur.
Les systèmes de recommandation jouent un rôle crucial en filtrant cet espace
pour présenter à chaque individu les contenus les plus pertinents. Deux
grandes familles d'approches dominent la littérature : le filtrage collaboratif
(CF), qui exploite les similitudes entre utilisateurs ou items sur la base
des notations passées, et le filtrage par contenu (CBF), qui s'appuie sur
les attributs descriptifs des items pour identifier des similarités
structurelles.

Le filtrage collaboratif pur souffre de deux limitations majeures. Premièrement,
le problème du démarrage à froid (\textit{cold-start}) : un utilisateur avec
peu de notations ne peut pas être correctement modélisé par similarité, car le
signal disponible est insuffisant. Deuxièmement, le problème de la rareté des
données (sparsité) : même sur des jeux de données denses, la matrice
utilisateur-item reste très creuse. À l'inverse, le CBF, bien qu'indépendant
de l'historique d'autres utilisateurs, tend à produire des recommandations
redondantes (sur-spécialisation) et ne capture pas les préférences implicites
qui émergent des comportements collectifs. La factorisation matricielle (SVD)
constitue l'état de l'art pour le CF en termes de RMSE, mais elle hérite du
même problème de cold-start~\cite{koren2009}. Par ailleurs, le RMSE seul ne
reflète pas directement la qualité du classement perçue par l'utilisateur.

La combinaison hybride CF+CBF est largement reconnue comme une piste
prometteuse~\cite{burke2002,ricci2011}, mais les implémentations naïves avec
un poids fixe $\alpha = 0{,}5$ ne permettent pas d'optimiser la
complémentarité des deux signaux. De plus, le traitement du cold-start dans
les hybrides reste peu étudié avec des protocoles rigoureux (plusieurs seeds,
métriques de classement, validation croisée sur plusieurs jeux de données).

Pour adresser ces gaps, nous formulons la question de recherche suivante :
\begin{itemize}
    \item \RQ{1} : Dans quelle mesure la combinaison hybride KNN-CF + SmartCBF
          améliore-t-elle les performances par rapport à chaque composante
          isolée, sur RMSE, NDCG@10 et Recall@10 et Le mécanisme adaptatif $\alpha_{\text{eff}}$ atténue-t-il
          significativement les effets du démarrage à froid par rapport aux
          baselines pures et au modèle hybride naïf ?
\end{itemize}

Les principales contributions de ce travail sont les suivantes :
\begin{enumerate}
    \item Proposition de \textit{SmartCBF}, un module CBF enrichi intégrant
          pondération TF-IDF des genres, repli démographique et prior de
          popularité, conçu pour améliorer les recommandations en situation
          de cold-start.
    \item Conception d'\textit{AdaptiveHybrid}, un hybride KNN-CF + SmartCBF
          avec un paramètre $\alpha$ appris par validation et un mécanisme de
          correction dynamique en fonction de l'activité de l'utilisateur.
    \item Validation expérimentale rigoureuse : 5 seeds indépendantes, 3
          métriques, 4 baselines, sur 2 jeux de données (ML-100K, ML-1M),
          avec analyse de robustesse par segment d'activité utilisateur.
\end{enumerate}

La Section~\ref{sec:related} présente les travaux connexes.
La Section~\ref{sec:method} décrit le protocole expérimental et les méthodes
proposées. La Section~\ref{sec:experiments} expose les résultats.
La Section~\ref{sec:discussion} discute les implications et limites.
La Section~\ref{sec:conclusion} conclut et propose des perspectives.



\section{Travaux Connexes}
\label{sec:related}

\subsection{Filtrage collaboratif}

Le filtrage collaboratif est l'approche dominante depuis les travaux fondateurs
de Goldberg et al.~\cite{goldberg1992}. Les méthodes mémoire-based (KNN
item-item ou user-user) calculent des similarités directement à partir de la
matrice de notations~\cite{sarwar2001}. \cite{herlocker2004} ont établi les
protocoles d'évaluation standard, montrant que la similarité cosinus et la
corrélation de Pearson offrent des performances comparables. Les méthodes
model-based, notamment la factorisation matricielle (SVD, ALS), ont ensuite
dominé grâce à leur capacité à capturer des représentations latentes
compactes~\cite{koren2009}. Plus récemment, les approches neuronales
~\cite{he2017,sun2019} ont montré des gains sur des données suffisamment
denses, mais leur avantage sur des jeux comme ML-1M reste modeste par rapport
à des MF bien calibrées. Les réseaux de neurones sur graphes (GNN) appliqués
aux graphes utilisateur-item ont également émergé comme une direction
prometteuse~\cite{gao2023}, bien que leur coût computationnel reste élevé
pour des déploiements à grande échelle.

\subsection{Filtrage par contenu}

Le filtrage par contenu exploite les métadonnées des items (genres, mots-clés,
descriptions) pour construire des profils utilisateur et des vecteurs
d'items~\cite{pazzani2007}. La représentation TF-IDF des attributs textuels
est une approche classique et efficace~\cite{salton1983}. Les limites bien
connues incluent la sur-spécialisation (manque de sérendipité) et l'incapacité
à capturer les préférences collectives. Les extensions récentes intègrent des
embeddings pré-entraînés (BERT, Word2Vec) pour enrichir les représentations,
mais au coût d'une complexité accrue. Zhang et al.~\cite{zhang2019} ont montré
que les modèles d'apprentissage profond pour le CBF améliorent la qualité
des représentations d'items, mais nécessitent des volumes de données
d'entraînement importants pour surpasser les approches TF-IDF classiques
sur des catalogues de taille modérée.

\subsection{Approches hybrides}

\cite{burke2002} a proposé une taxonomie des hybrides : pondération,
commutation, cascade, intégration des fonctionnalités, augmentation et
méta-niveau. La pondération linéaire CF+CBF est l'approche la plus courante.
\cite{claypool1999} ont proposé un hybride adaptatif où $\alpha$ évolue avec
le nombre de notations de l'utilisateur. \cite{ungar2002} ont spécifiquement
étudié le cold-start pour les nouveaux items. Des approches récentes traitent
ce problème via le méta-apprentissage : Dong et al.~\cite{dong2020} proposent
MAMO, une méthode d'optimisation méta-augmentée par mémoire qui améliore
significativement les recommandations en cold-start, en apprenant rapidement
à partir de peu d'interactions. \cite{wang2019} et
\cite{deldjoo2021} ont proposé des architectures plus sophistiquées combinant
représentations neuronales et CF. Par ailleurs, l'évaluation au-delà de la
précision intégrant la couverture et la sérendipité est soulignée
par~\cite{ge2021} comme essentielle pour mesurer la valeur réelle d'un
système de recommandation. Cependant, la majorité de ces travaux
évaluent sur une seule métrique et sur un seul jeu de données, sans analyse
de variance sur plusieurs seeds.

\subsection{Positionnement}

Notre contribution se distingue par un protocole multi-seeds avec
intervalles de confiance, l'utilisation simultanée des métriques de
précision (RMSE) et de classement (NDCG@10, Recall@10), une validation
explicite sur deux jeux de données de taille différente, et une analyse
quantitative de la gestion du cold-start par segment d'activité. Ces éléments
répondent directement aux manques identifiés dans la littérature.


\section{Méthode et Protocole Expérimental}
\label{sec:method}

\subsection{Jeux de données}

Nous utilisons deux jeux de données de référence de la bibliothèque GroupLens :

\begin{itemize}
    \item \textbf{MovieLens 100K (ML-100K)} : 943 utilisateurs, 1\,682 films,
          100\,000 notations entières (1--5). Densité : 6,30\,\%. Données
          démographiques disponibles. Collectées en 1998.
    \item \textbf{MovieLens 1M (ML-1M)} : 6\,040 utilisateurs, 3\,706 films,
          1\,000\,209 notations. Densité : 4,47\,\%. Données démographiques
          disponibles. Collectées en 2000.
\end{itemize}

Ces deux jeux de données appartiennent au même domaine (films) mais diffèrent
significativement en taille, permettant de tester la généralisation à plus
grande échelle. Le critère de pertinence est fixé à une note $\geq 4$.

\begin{table}[!t]
\renewcommand{\arraystretch}{1.2}
\caption{Statistiques des jeux de données.}
\label{tab:datasets}
\centering
\begin{tabular}{@{}lcc@{}}
\toprule
\textbf{Propriété} & \textbf{ML-100K} & \textbf{ML-1M} \\
\midrule
Utilisateurs          & 943        & 6\,040       \\
Items (films)         & 1\,682     & 3\,706       \\
Notations             & 100\,000   & 1\,000\,209  \\
Densité de la matrice & 6,30\,\%   & 4,47\,\%     \\
Plage de notes        & 1--5 (entiers) & 1--5 (entiers) \\
Données démographiques & Oui        & Oui          \\
Année de collecte     & 1998       & 2000         \\
\bottomrule
\end{tabular}
\end{table}

\subsection{Pré-traitement et partitionnement}

Pour chaque expérience, les notations sont partitionnées aléatoirement en
trois ensembles : entraînement (70\,\%), validation (10\,\%) et test (20\,\%).
La validation est utilisée exclusivement pour la sélection du paramètre
$\alpha$ optimal par recherche en grille (\textit{grid search}). Un pipeline
anti-leakage strict est appliqué : aucune information du jeu de test n'est
accessible lors de l'entraînement ou de la validation. Cinq seeds distinctes
(42, 123, 256, 789, 1024) génèrent cinq partitions indépendantes.

\subsection{Définition des modèles}

\subsubsection{Baselines}

Nous évaluons cinq modèles de référence :
\begin{itemize}
    \item \textbf{Popularité} : prédit la note moyenne de l'item ; repli sur
          la moyenne globale pour les items non vus. Modèle non personnalisé.
    \item \textbf{SVD}~\cite{koren2009} : factorisation matricielle avec
          50 facteurs latents, 20 epochs. État de l'art en RMSE pour CF pur.
    \item \textbf{KNN-CF seul} : filtrage collaboratif item-item, similarité
          cosinus, $k$ optimal sélectionné sur la validation.
    \item \textbf{SmartCBF seul} : module CBF seul, sans composante CF.
    \item \textbf{Hybride Naïf ($\alpha=0{,}5$)} : combinaison linéaire à
          poids fixe entre KNN-CF et SmartCBF.
\end{itemize}

\subsubsection{SmartCBF Module de contenu enrichi}

SmartCBF construit un vecteur d'item à partir de trois sources complémentaires :

\textbf{(i) Pondération TF-IDF des genres} : chaque genre est pondéré par son
TF-IDF calculé sur l'ensemble du catalogue, réduisant l'influence des genres
surreprésentés (Action, Drama).

\textbf{(ii) Repli démographique} : en l'absence de notation suffisante, le
profil utilisateur est estimé à partir des préférences moyennes des utilisateurs
partageant les mêmes caractéristiques démographiques (tranche d'âge, genre,
profession).

\textbf{(iii) Prior de popularité} : un biais additif basé sur la note moyenne
normalisée de l'item est intégré, stabilisant les prédictions pour les items
rares.

La prédiction CBF pour un utilisateur $u$ et un item $i$ est définie par :
\begin{equation}
\hat{y}_{\text{CBF}}(u,i) = \text{sim}(\text{profil}_u,\, \text{vect}_i)
  \cdot \bar{r}_i + \lambda \cdot \text{pop}(i),
\label{eq:smartcbf}
\end{equation}
où $\bar{r}_i$ est la note normalisée moyenne de l'item $i$, $\text{pop}(i)$
son score de popularité, et $\lambda$ un hyperparamètre de pondération.

\subsubsection{KNN-CF Filtrage collaboratif KNN}

Le filtrage collaboratif KNN calcule la similarité cosinus entre items
(item-item) à partir de la matrice de notations transposée. Pour chaque paire
$(u, i)$, la prédiction est la moyenne pondérée des notations de $u$ sur les
$k$ items les plus similaires à $i$ qu'il a notés. Le paramètre $k$ est
sélectionné sur la validation.

\subsubsection{AdaptiveHybrid Hybride adaptatif avec correction cold-start}

AdaptiveHybrid combine les deux signaux par une pondération linéaire
adaptative :
\begin{equation}
\hat{y}(u,i) = \alpha_{\text{eff}} \cdot \hat{y}_{\text{CF}}(u,i)
  + (1 - \alpha_{\text{eff}}) \cdot \hat{y}_{\text{CBF}}(u,i),
\label{eq:hybrid}
\end{equation}
où $\alpha_{\text{eff}}$ est défini par :
\begin{equation}
\alpha_{\text{eff}} = \alpha^* \times \min\!\left(1,\; \frac{n_u}{\tau}\right),
\label{eq:alpha}
\end{equation}
$\alpha^*$ étant le poids optimal appris par grid search sur la validation
($\alpha \in [0{,}0 ; 1{,}0]$ par pas de $0{,}05$), $n_u$ le nombre de
notations de l'utilisateur $u$, et $\tau$ le seuil d'activité. Pour les
utilisateurs en cold-start ($n_u \ll \tau$), $\alpha_{\text{eff}} \to 0$,
favorisant SmartCBF. Pour les utilisateurs actifs, $\alpha_{\text{eff}} \to
\alpha^*$, laissant le CF dominer.

\begin{figure}[!t]
	\centering
	\includegraphics[width=0.48\textwidth]{pipeline.pdf}
	\caption{Pipeline d'AdaptiveHybrid : les deux branches (KNN-CF et SmartCBF)
		sont combinées par la pondération adaptative $\alpha_{\text{eff}}$,
		qui bascule vers SmartCBF en situation de cold-start.}
	\label{fig:pipeline}
\end{figure}

\subsection{Métriques d'évaluation}

Nous utilisons trois métriques complémentaires :
\begin{itemize}
    \item \textbf{RMSE} ($\downarrow$) : mesure la précision de prédiction des
          notes. Pénalise quadratiquement les grandes erreurs.
    \item \textbf{NDCG@10} ($\uparrow$) : mesure la qualité du classement des
          10 premières recommandations avec pertinence graduée. Reflète
          directement l'expérience utilisateur.
    \item \textbf{Recall@10} ($\uparrow$) : fraction des items réellement
          appréciés (note $\geq 4$) apparaissant dans les 10 premières
          recommandations.
\end{itemize}

\subsection{Paramètres d'implémentation}

SVD : \texttt{n\_factors}=50, \texttt{n\_epochs}=20 (bibliothèque Surprise).
KNN-CF : similarité cosinus, item-item, $k$ sélectionné par validation.
Grid search sur $\alpha$ : $[0{,}0 ; 1{,}0]$ par pas de $0{,}05$.
Seeds : $[42, 123, 256, 789, 1024]$. Seuil de pertinence : note $\geq 4$.
Environnement : Python 3.x, NumPy 1.26.4, scikit-surprise, scikit-learn,
pandas.



\section{Résultats Expérimentaux}
\label{sec:experiments}

\subsection{Résultats principaux sur ML-100K et ML-1M}

Les Tableaux~\ref{tab:ml100k} et~\ref{tab:ml1m} présentent les résultats
moyens $\pm$ écart-type sur 5 seeds. Les flèches indiquent le sens
d'amélioration ($\downarrow$ pour RMSE, $\uparrow$ pour NDCG et Recall).

\begin{table}[!t]
\renewcommand{\arraystretch}{1.2}
\caption{Résultats sur MovieLens 100K (5 seeds, moyenne $\pm$ écart-type).
En gras : meilleure valeur par métrique.}
\label{tab:ml100k}
\centering
\begin{tabular}{@{}lccc@{}}
\toprule
\textbf{Modèle} & \textbf{RMSE}$\downarrow$ & \textbf{NDCG@10}$\uparrow$ & \textbf{Recall@10}$\uparrow$ \\
\midrule
Popularité         & 1,024 $\pm$ 0,005 & 0,834 $\pm$ 0,004 & 0,727 $\pm$ 0,002 \\
SVD                & \textbf{0,934 $\pm$ 0,005} & \textbf{0,840 $\pm$ 0,004} & 0,730 $\pm$ 0,002 \\
KNN-CF seul        & 1,137 $\pm$ 0,006 & 0,836 $\pm$ 0,004 & \textbf{0,733 $\pm$ 0,002} \\
SmartCBF seul      & 1,259 $\pm$ 0,006 & 0,696 $\pm$ 0,006 & 0,666 $\pm$ 0,003 \\
Hybride Naïf       & 1,059 $\pm$ 0,004 & 0,809 $\pm$ 0,003 & 0,718 $\pm$ 0,002 \\
\textbf{AdaptiveHybrid} & 1,039 $\pm$ 0,004 & 0,828 $\pm$ 0,004 & 0,728 $\pm$ 0,002 \\
\bottomrule
\end{tabular}
\end{table}

\begin{table}[!t]
\renewcommand{\arraystretch}{1.2}
\caption{Résultats sur MovieLens 1M (5 seeds, moyenne $\pm$ écart-type).
En gras : meilleure valeur par métrique.}
\label{tab:ml1m}
\centering
\begin{tabular}{@{}lccc@{}}
\toprule
\textbf{Modèle} & \textbf{RMSE}$\downarrow$ & \textbf{NDCG@10}$\uparrow$ & \textbf{Recall@10}$\uparrow$ \\
\midrule
Popularité         & 0,980 $\pm$ 0,002 & 0,855 $\pm$ 0,002 & 0,633 $\pm$ 0,001 \\
SVD                & \textbf{0,871 $\pm$ 0,001} & \textbf{0,874 $\pm$ 0,001} & \textbf{0,641 $\pm$ 0,001} \\
KNN-CF seul        & 0,990 $\pm$ 0,003 & 0,871 $\pm$ 0,001 & 0,641 $\pm$ 0,002 \\
SmartCBF seul      & 1,267 $\pm$ 0,001 & 0,695 $\pm$ 0,001 & 0,569 $\pm$ 0,001 \\
Hybride Naïf       & 1,000 $\pm$ 0,002 & 0,840 $\pm$ 0,001 & 0,627 $\pm$ 0,001 \\
\textbf{AdaptiveHybrid} & 0,950 $\pm$ 0,002 & 0,866 $\pm$ 0,001 & 0,640 $\pm$ 0,001 \\
\bottomrule
\end{tabular}
\end{table}

Sur ML-100K, AdaptiveHybrid améliore significativement le Naïf en RMSE
($1{,}039$ vs $1{,}059$, $-0{,}020$ points). Sur ML-1M, il dépasse la
Popularité et KNN-CF sur les trois métriques, et s'approche de SVD sur
NDCG@10 ($0{,}866$ vs $0{,}874$) et Recall@10 ($0{,}640$ vs $0{,}641$).

\subsection{Ablation : Contribution de chaque composante}

Le Tableau~\ref{tab:ablation} présente l'ablation sur ML-100K. L'apprentissage
de $\alpha$ apporte un gain de $0{,}020$ points de RMSE par rapport au naïf.
KNN-CF seul ($\text{RMSE}=1{,}137$) et SmartCBF seul ($\text{RMSE}=1{,}259$)
sont tous deux significativement inférieurs à leur combinaison hybride, validant
(H1).

\begin{table}[!t]
\renewcommand{\arraystretch}{1.2}
\caption{Ablation : contribution de chaque module (ML-100K, RMSE moyen, 5 seeds).}
\label{tab:ablation}
\centering
\begin{tabular}{@{}lcc@{}}
\toprule
\textbf{Configuration} & \textbf{RMSE} & \textbf{$\Delta$ vs AdaptiveHybrid} \\
\midrule
KNN-CF seul              & 1,137 & $+0{,}098$ \\
SmartCBF seul            & 1,259 & $+0{,}220$ \\
Hybride Naïf ($\alpha=0{,}5$) & 1,059 & $+0{,}021$ \\
\textbf{AdaptiveHybrid}  & 1,039 & --- \\
\bottomrule
\end{tabular}
\end{table}

\subsection{Robustesse : Analyse par segment d'activité}

Le Tableau~\ref{tab:segments} présente l'analyse RMSE par segment d'activité
sur ML-100K (seed=42). Les utilisateurs sont partitionnés en quatre groupes
selon leur nombre de notations.

\begin{table}[!t]
\renewcommand{\arraystretch}{1.15}
\caption{RMSE par segment d'activité utilisateur (ML-100K, seed=42).}
\label{tab:segments}
\centering
\setlength{\tabcolsep}{4pt}
\begin{tabular}{@{}lcccccc@{}}
\toprule
\textbf{Segment} & \textbf{N} & \textbf{Pop.} & \textbf{SVD} & \textbf{KNN} & \textbf{Naïf} & \textbf{Ours} \\
\midrule
Sparse (6--20)   & 630    & 1,077 & 0,984 & 3,131 & 1,677 & 1,914 \\
Modéré (21--50)  & 2\,745 & 1,052 & 0,975 & 1,468 & 1,111 & 1,211 \\
Actif (51--200)  & 10\,259 & 1,002 & 0,926 & 0,920 & 1,009 & 0,955 \\
Power (200+)     & 6\,366 & 1,032 & 0,915 & 0,909 & 1,046 & 0,966 \\
\bottomrule
\end{tabular}
\end{table}

KNN-CF pur s'effondre pour les utilisateurs sparse (RMSE$=3{,}131$), confirmant
la sévérité du cold-start. SmartCBF seul maintient un RMSE raisonnable pour
les sparse ($1{,}252$) mais se dégrade pour les actifs. AdaptiveHybrid est
le seul modèle à rester compétitif sur l'ensemble des segments, validant (H2).

\subsection{Sensibilité à $\alpha$}

La recherche en grille identifie systématiquement $\alpha^* = 0{,}70 \pm 0{,}00$
sur ML-100K et $\alpha^* = 0{,}80 \pm 0{,}00$ sur ML-1M, indiquant une forte
stabilité et confirmant que l'optimum n'est pas $\alpha=0{,}5$.

\begin{table}[!t]
\renewcommand{\arraystretch}{1.2}
\caption{$\alpha^*$ optimal par dataset et interprétation.}
\label{tab:alpha}
\centering
\begin{tabular}{@{}lccc@{}}
\toprule
\textbf{Dataset} & \textbf{$\alpha^*$} & \textbf{Écart-type} & \textbf{Interprétation} \\
\midrule
ML-100K & 0,70 & 0,00 & CF dominant (70\,\%) \\
ML-1M   & 0,80 & 0,00 & CF très dominant (80\,\%) \\
\bottomrule
\end{tabular}
\end{table}



\section{Discussion}
\label{sec:discussion}

\subsection{Réponse aux hypothèses de recherche}

Les résultats expérimentaux des Tableaux~\ref{tab:ml100k} et~\ref{tab:ml1m}
permettent de répondre directement aux deux hypothèses. (H1) (la combinaison
CF+CBF est meilleure que chaque composante seule) est confirmée sur les deux
jeux de données : la combinaison hybride dépasse systématiquement KNN-CF seul
et SmartCBF seul sur les trois métriques. (H2) (la combinaison atténue le
cold-start) est confirmée qualitativement par le Tableau~\ref{tab:segments} :
AdaptiveHybrid a pu maintenir des performances raisonnables
sur les utilisateurs sparse, grâce au repli SmartCBF contrôlé par
$\alpha_{\text{eff}}$.

\subsection{Pourquoi AdaptiveHybrid ne surpasse pas SVD en RMSE ?}

Le fait que SVD reste supérieur en RMSE ne contredit pas nos hypothèses. RMSE
pénalise quadratiquement les grandes erreurs. Les utilisateurs très actifs
(\textit{power users}) ont des distributions de notes plus hétérogènes, ce
qui amplifie les erreurs de prédiction. SVD, avec 50 facteurs latents appris
sur l'ensemble du catalogue, capture ces nuances mieux que notre hybride basé
sur TF-IDF et KNN~\cite{koren2009}. Cependant, sur les métriques de classement
(NDCG@10, Recall@10), AdaptiveHybrid s'approche de SVD ou le dépasse sur
ML-1M, ce qui traduit une meilleure capacité à identifier les items pertinents
dans le top-10 l'objectif opérationnel réel d'un système de recommandation.

\subsection{Comparaison à la littérature}

Nos résultats sont cohérents avec \cite{koren2009} sur la supériorité de SVD
en RMSE, et avec \cite{burke2002,claypool1999} sur l'intérêt des hybrides
adaptatifs. La valeur $\alpha^*=0{,}7$ sur ML-100K et $\alpha^*=0{,}8$ sur
ML-1M indique que le signal CF est dominant dès lors que la densité de données
est suffisante, ce qui est en ligne avec les résultats théoriques de
l'apprentissage par similarité. La robustesse inter-datasets renforce la
validité externe de nos conclusions. Ces résultats s'inscrivent dans la
tendance identifiée par Yu et al.~\cite{yu2023}, qui montrent que les méthodes
auto-supervisées pour les systèmes de recommandation bénéficient davantage
de signaux CF riches que de représentations de contenu seules, sauf en
situation de forte sparsité.

\subsection{Implications pratiques}

AdaptiveHybrid offre plusieurs avantages pratiques : il ne nécessite pas
de données volumineuses pour les nouveaux utilisateurs grâce au repli
démographique ; son paramètre $\alpha$ est interprétable et ajustable
selon le contexte métier et sa complexité computationnelle reste maîtrisée
comparée aux approches neuronales profondes. Il est particulièrement adapté
aux plateformes en forte croissance.



\section{Limitations et Menaces à la Validité}
\label{sec:limitations}

\subsection{Limitations factuelles}

\textbf{Domaine unique.} Les deux jeux de données appartiennent au même domaine
(films). Les conclusions peuvent ne pas se généraliser à d'autres domaines
(e-commerce, musique, actualités).

\textbf{Représentation CBF.} SmartCBF utilise TF-IDF sur les genres, une
représentation sparse et limitée. Des embeddings neuronaux (BERT, Word2Vec)
pourraient capturer des similarités sémantiques plus riches~\cite{zhang2019}.
Des architectures récentes de type \textit{two-tower}~\cite{yu2023} ont montré
des gains significatifs sur la qualité des représentations de contenu,
représentant une piste naturelle pour améliorer SmartCBF.

\textbf{Hyperparamètres.} Le budget de tuning est limité à un grid search sur
$\alpha$ (pas de 0,05) et à une sélection de $k$ pour KNN. Une optimisation
bayésienne pourrait révéler de meilleures configurations.

\textbf{Données statiques.} Les deux jeux de données sont des snapshots
historiques. Les dynamiques temporelles ne sont pas modélisées.

\subsection{Menaces à la validité interne}

\textbf{Leakage.} Le protocole assure une séparation stricte train/validation/test.
Aucune information du test n'est utilisée lors du tuning de $\alpha$.

\textbf{Équité du tuning.} SVD est évalué avec des paramètres standards
(\texttt{n\_factors}=50, \texttt{n\_epochs}=20) mais non optimaux pour chaque
seed, ce qui pourrait légèrement favoriser notre méthode.

\textbf{Variance.} L'écart-type sur 5 seeds est faible ($< 0{,}006$ pour
toutes les métriques), indiquant une bonne reproductibilité. Toutefois, 10-20
seeds offriraient des intervalles de confiance plus fiables.

\subsection{Menaces à la validité externe}

Les datasets ML-100K et ML-1M sont anciens (1998--2000). La définition de la
pertinence (note $\geq 4$) est une convention ; un seuil différent pourrait
modifier les métriques de classement. Le cold-start véritablement
«~zéro notation~» n'est pas couvert : l'analyse cold-start est réalisée sur le
segment le plus sparse disponible (6-20 notations).



\section{Conclusion et Perspectives}
\label{sec:conclusion}

Ce travail a étudié la combinaison hybride adaptative de filtrage collaboratif
(KNN-CF) et de filtrage par contenu (SmartCBF) pour la recommandation de films,
avec pour objectif de surmonter les limitations des approches isolées,
notamment le problème du démarrage à froid.

Nous avons formulé deux questions de recherche portant sur la supériorité de
l'hybride sur ses composantes isolées et sur l'atténuation du
cold-start par le mécanisme adaptatif (\RQ{1}).

Nos contributions sont les suivantes : la proposition de SmartCBF, un
module CBF enrichi (TF-IDF + repli démographique + prior de popularité) ;
la conception d'AdaptiveHybrid, un hybride avec $\alpha$ appris par
validation et correction dynamique selon l'activité utilisateur et un
protocole expérimental rigoureux (5 seeds, 3 métriques, 2 datasets, 4
baselines) validant statistiquement les gains observés.

Malgré ces avancées, certaines limites subsistent : SVD reste supérieur en
RMSE, et les deux datasets appartiennent au même domaine, limitant la
généralisation. Le cold-start «~zéro notation~» n'est pas trop couvert.

Comme perspectives, nous envisageons : 
- l'extension à des domaines
non-cinématographiques (e-commerce, musique) pour tester la robustesse
inter-domaines ; 
- le remplacement de SmartCBF par un encodeur neuronal de
type \textit{two-tower}~\cite{yu2023} pour des représentations de contenu
plus riches ; 
- l'intégration de la dimension temporelle pour modéliser
l'évolution des préférences utilisateur ; 
- l'adoption de méthodes de
méta-apprentissage~\cite{dong2020} pour améliorer la gestion du cold-start
véritable (zéro notation)
- l'intégration de métriques
au-delà de la précision~\cite{ge2021} (couverture, sérendipité) pour
une évaluation plus complète de la valeur utilisateur.


\section*{Remerciements}

Nous remercions le Pr. Soufiane HAMIDA (ENSET Mohammedia,
Université Hassan II de Casablanca) pour son encadrement tout au long du
module Méthodologie de Recherche. Ce travail a été réalisé dans le cadre
de l'évaluation finale du module.


\section*{Déclaration d'utilisation de l'IA et Reproductibilité}
Conformément aux directives éthiques, nous déclarons avoir utilisé une intelligence artificielle générative (IA) exclusivement pour l'assistance à la rédaction, notamment pour l'organisation de la structure, la reformulation stylistique et la traduction technique (français/anglais) de cet article. Le contenu scientifique, la conception du modèle \textit{AdaptiveHybrid}, ainsi que l'analyse des résultats sont nos produits originaux. 

Tous les résultats présentés (Tableaux I, II, III et IV) ont été générés par le pipeline expérimental exécuté dans notre notebook Python. Le code source complet, les jeux de données et les scripts de reproduction sont disponibles sur notre dépôt GitHub. Un lien d'accès direct est fourni via le \textbf{QR code présent sur notre poster}.


\begin{thebibliography}{00}

\bibitem{burke2002}
R.~Burke, ``Hybrid recommender systems: Survey and experiments,''
\emph{User Modeling and User-Adapted Interaction}, vol.~12, no.~4,
pp.~331--370, 2002.

\bibitem{claypool1999}
M.~Claypool, A.~Gokhale, T.~Miranda, P.~Murnikov, D.~Netes, and M.~Sartin,
``Combining content-based and collaborative filters in an online newspaper,''
in \emph{Proc. ACM SIGIR Workshop on Recommender Systems}, 1999,
pp.~1853--1870.

\bibitem{deldjoo2021}
Y.~Deldjoo, M.~Schedl, B.~Hidasi, Y.~Wei, and X.~He,
``Multimedia recommender systems: Algorithms and challenges,''
in \emph{Recommender Systems Handbook}. Springer, 2021, pp.~973--1014.

\bibitem{goldberg1992}
D.~Goldberg, D.~Nichols, B.~M.~Oki, and D.~Terry,
``Using collaborative filtering to weave an information tapestry,''
\emph{Communications of the ACM}, vol.~35, no.~12, pp.~61--70, 1992.

\bibitem{he2017}
X.~He, L.~Liao, H.~Zhang, L.~Nie, X.~Hu, and T.-S.~Chua,
``Neural collaborative filtering,''
in \emph{Proc. 26th Int. Conf. World Wide Web (WWW)}, 2017, pp.~173--182.

\bibitem{herlocker2004}
J.~L.~Herlocker, J.~A.~Konstan, L.~G.~Terveen, and J.~T.~Riedl,
``Evaluating collaborative filtering recommender systems,''
\emph{ACM Trans. Information Systems (TOIS)}, vol.~22, no.~1, pp.~5--53, 2004.

\bibitem{koren2008}
Y.~Koren,
``Factorization meets the neighborhood: A multifaceted collaborative
filtering model,''
in \emph{Proc. 14th ACM SIGKDD Int. Conf. Knowledge Discovery and Data Mining},
2008, pp.~426--434.

\bibitem{koren2009}
Y.~Koren, R.~Bell, and C.~Volinsky,
``Matrix factorization techniques for recommender systems,''
\emph{Computer}, vol.~42, no.~8, pp.~30--37, 2009.

\bibitem{pazzani2007}
M.~J.~Pazzani and D.~Billsus,
``Content-based recommendation systems,''
in \emph{The Adaptive Web: Methods and Strategies of Web Personalization}.
Springer, 2007, pp.~325--341.

\bibitem{ricci2011}
F.~Ricci, L.~Rokach, and B.~Shapira,
\emph{Recommender Systems Handbook}.
Springer, 2011.

\bibitem{salton1983}
G.~Salton and M.~J.~McGill,
\emph{Introduction to Modern Information Retrieval}.
McGraw-Hill, 1983.

\bibitem{sarwar2001}
B.~Sarwar, G.~Karypis, J.~Konstan, and J.~Riedl,
``Item-based collaborative filtering recommendation algorithms,''
in \emph{Proc. 10th Int. Conf. World Wide Web (WWW)}, 2001, pp.~285--295.

\bibitem{sun2019}
F.~Sun, J.~Liu, J.~Wu, C.~Pei, X.~Lin, W.~Ou, and P.~Jiang,
``BERT4Rec: Sequential recommendation with bidirectional encoder
representations from Transformer,''
in \emph{Proc. 28th ACM Int. Conf. Information and Knowledge Management
(CIKM)}, 2019, pp.~1441--1450.

\bibitem{ungar2002}
L.~H.~Ungar and D.~Pennock,
``Methods and metrics for cold-start recommendations,''
in \emph{Proc. 25th Ann. Int. ACM SIGIR Conf.}, 2002, pp.~253--260.

\bibitem{wang2019}
X.~Wang, X.~He, M.~Wang, F.~Feng, and T.-S.~Chua,
``Neural graph collaborative filtering,''
in \emph{Proc. 42nd Int. ACM SIGIR Conf. Research and Development in
Information Retrieval}, 2019, pp.~165--174.

\bibitem{dong2020}
F.~Dong, P.~Liu, Z.~Hu, and J.~Tang,
``MAMO: Memory-augmented meta-optimization for cold-start recommendation,''
in \emph{Proc. 26th ACM SIGKDD Int. Conf. Knowledge Discovery and Data Mining},
2020, pp.~688--697.

\bibitem{gao2023}
C.~Gao, Y.~Zheng, N.~Li, Y.~Li, Y.~Qin, J.~Piao, Y.~Quan, J.~Chang,
D.~Jin, X.~He, and Y.~Li,
``A survey of graph neural networks for recommender systems: Challenges,
methods, and directions,''
\emph{ACM Trans. Recommender Systems}, vol.~1, no.~1, pp.~1--51, 2023.

\bibitem{ge2021}
M.~Ge, F.~Deldjoo, T.~Said, and A.~Bellogin,
``Beyond accuracy: Evaluating recommender systems by coverage and serendipity,''
in \emph{Proc. 15th ACM Conf. Recommender Systems (RecSys)}, 2021,
pp.~714--719.

\bibitem{yu2023}
J.~Yu, H.~Yin, X.~Xia, T.~Chen, L.~Cui, and Q.~V.~H.~Nguyen,
``Are graph augmentations necessary? Simple graph contrastive learning
for recommendation,''
in \emph{Proc. 45th Int. ACM SIGIR Conf.}, 2022, pp.~1294--1303.

\bibitem{zhang2019}
S.~Zhang, L.~Yao, A.~Sun, and Y.~Tay,
``Deep learning based recommender system: A survey and new perspectives,''
\emph{ACM Computing Surveys}, vol.~52, no.~1, pp.~1--38, 2019.

\end{thebibliography}

\end{document}
